\chapter{Primitive operations: core}
\label{chap:primitives-core}

\section*{At a glance}
\begin{itemize}
\item \lead{Goal} Verify the kernel's recognition of GMP-accelerated primitives, short of the
eighteen per-primitive branches (Chapter~\ref{chap:primitives-arith}).
\item \lead{Where} \texttt{Primitive.checkDef} (\srcloc{Lean4Lean/Primitive.lean}{561}) matches a
name and type syntactically, then runs \texttt{isDefEq} on the definition's \hl{value}: open
equations such as $f\,x\,0 \equiv x$. This chapter closes those into $\mathsf{Reflects}$
statements and transports them onto the constant. Notation: $\Gamma \vdash e : A$ is
\texttt{VEnv.HasType}, $\equiv$ is the type-erased \texttt{VEnv.IsDefEqU}, $\overline{k}$ is
\texttt{VExpr.natLit k}, and a universe count $U$ is omitted when $0$.
\item \lead{Main results} Theorem~\ref{vpc:thm:reflects-toconst} carries a value's reflection onto
its constant; Theorem~\ref{vpc:thm:unfoldnatwf} discharges well-founded recursion;
Theorem~\ref{vpc:thm:checkdef-wf} is the boundary the declaration checker consumes.
\item \lead{Status} \stProved{} the closed vocabulary with no \texttt{axfoot} (telescopes,
\texttt{ArgsTyped}, \texttt{addDefEq}'s equation); \stTainted{} the reflection schemas and every
node past \ref{vpc:thm:prim-trnat-trbool}, through the admitted \texttt{VEnv.WF.patsStrong} and
through \texttt{IsDefEq.uniq}, including \texttt{checkDef.WF} itself and every
\texttt{PrimitiveResult}.
\item \lead{Lean files} \texttt{Lean4Lean/Verify/Primitive.lean},
\texttt{Lean4Lean/Verify/Environment/Primitive/Basic.lean},
\texttt{Lean4Lean/Verify/Environment/Primitive/Recursion.lean},
\texttt{Lean4Lean/Verify/Environment/Primitive.lean}.
\item \lead{Thesis} No counterpart: Carneiro's system has no GMP-accelerated primitives.
Individual nodes cite Weakening, Regularity, definitional inversion, and the beta/delta rules
where they ground a step; a node with no \texttt{Thesis} line has none.
\end{itemize}

\par\smallskip\noindent
\begin{tabular}{p{0.28\linewidth}p{0.52\linewidth}l}
\toprule
Operation & Specification & Status\\
\midrule
\texttt{Nat.add}, \texttt{Nat.sub} & unary step, Theorem~\ref{vpc:thm:reflects-unary-step} & \stTainted{}\\
\texttt{Nat.mul}, \texttt{Nat.pow} & binary step, Theorem~\ref{vpc:thm:reflects-binary-step} & \stTainted{}\\
\texttt{Nat.shiftLeft} & first-argument step, Lemma~\ref{vpc:thm:reflects-first-arg-step} & \stTainted{}\\
\texttt{Nat.beq}, \texttt{Nat.ble} & four constructor cases, Lemma~\ref{vpc:thm:reflects-ctor-cases} & \stTainted{}\\
\texttt{Nat.land}, \texttt{Nat.lor}, \texttt{Nat.xor} & boolean operator cases, Lemma~\ref{vpc:thm:reflects-bool-cases} & \stTainted{}\\
\texttt{Nat.pred} & two probed equations, Lemma~\ref{vpc:thm:reflects-pred-eqns} & \stTainted{}\\
\texttt{Nat.div}, \texttt{Nat.mod} & fuel recursion, Theorem~\ref{vpc:thm:natfuelrec} & \stProved{}\\
\texttt{Nat.gcd} & well-founded recursion, Theorem~\ref{vpc:thm:unfoldnatwf} & \stProved{}\\
\texttt{Nat.bitwise} & well-founded recursion, higher-order obligation, Theorem~\ref{vpc:thm:data-mkresultbitwise} & \stTainted{}\\
\texttt{Char.ofNat} & recorded type only, Theorem~\ref{vpc:thm:data-mkresult} & \stProved{}\\
\bottomrule
\end{tabular}\par\smallskip

\section{Reflection for the primitive definitions}

This section covers \srcloc{Lean4Lean/Verify/Primitive.lean}{1}: closed-context typing facts, the
five recurrence schemas, and the transfer from value to constant.

\begin{lemma}[Numerals are closed]
\label{vpc:lem:prim-closedn-natlit}\lean{Lean4Lean.VExpr.closedN_natLit}\leanok
\uses{def:closedn}
\srcloc{Lean4Lean/Verify/Primitive.lean}{90}

\lead{In short} The model numeral $\overline{n}$ is \texttt{ClosedN k} at every de Bruijn depth
$k$.
\end{lemma}
\begin{proof}\leanok
\lead{Idea} Induction on $n$, unfolding \texttt{VExpr.natLit}, \texttt{natZero} and
\texttt{natSucc}.
\end{proof}

\begin{lemma}[Nat/Bool typing at any context]
\label{vpc:thm:prim-nat-bool-typing}
\lean{Lean4Lean.VEnv.HasPrimitives.natZeroT, Lean4Lean.VEnv.HasPrimitives.natSuccT, Lean4Lean.VEnv.HasPrimitives.natPredT, Lean4Lean.VEnv.HasPrimitives.natLitT, Lean4Lean.VEnv.HasPrimitives.boolLitT, Lean4Lean.VEnv.HasPrimitives.natIsType, Lean4Lean.VEnv.HasPrimitives.boolIsType}\leanok
\uses{family:hasprimitives-accessors}
\contrib{mixed}\srcloc{Lean4Lean/Verify/Primitive.lean}{29}
\thesisref{Weakening and Regularity, typesys.tex thm:weak, thm:reg.}

\lead{In short} Seven wrappers, assuming \texttt{env.OrderedStrong}, moving
\texttt{HasPrimitives}'s closed-context facts to an \hl{arbitrary context} $\Gamma$: the form
every proof below uses.
\begin{itemize}
\item $\Gamma \vdash \mathtt{Nat.zero} : \mathtt{nat}$, $\Gamma \vdash \mathtt{Nat.succ} :
\mathtt{nat} \to \mathtt{nat}$, $\Gamma \vdash \mathtt{Nat.pred} : \mathtt{nat} \to \mathtt{nat}$
(given containment), $\Gamma \vdash \overline{k} : \mathtt{nat}$ at any $U$, and $\Gamma \vdash
\mathtt{boolLit}\ b : \mathtt{bool}$;
\item $\mathtt{nat}$ and $\mathtt{bool}$ are types in any well-formed \texttt{VLCtx}.
\end{itemize}
\end{lemma}
\begin{proof}\leanok
\uses{family:trexprs-literals, thm:isdefeq-istype}
\lead{Idea} Each is the corresponding \texttt{TrExprS} literal rule at the empty context, pushed
to $\Gamma$ by \texttt{HasType.weak0}; the two \texttt{IsType} members take
\texttt{IsDefEq.isType} of the constructor typing instead.
\end{proof}

\begin{lemma}[Nat/Bool containment from a recorded type]
\label{vpc:thm:prim-contains-nat}
\lean{Lean4Lean.VEnv.contains_nat_of_hasType, Lean4Lean.VEnv.HasPrimitives.natOfPred, Lean4Lean.VEnv.HasPrimitives.natOfAdd, Lean4Lean.VEnv.HasPrimitives.natOfMul, Lean4Lean.VEnv.HasPrimitives.natOfDiv, Lean4Lean.VEnv.HasPrimitives.natOfMod, Lean4Lean.VEnv.HasPrimitives.boolOfBitwise}\leanok
\uses{family:hasprimitives-accessors}
\contrib{mixed}\srcloc{Lean4Lean/Verify/Primitive.lean}{47}
\thesisref{Regularity, definitional inversion; typesys.tex thm:reg, unique.tex thm:1dinv.}

\lead{In short} If $\vdash e : \mathtt{nat} \to A$ at the empty context, \texttt{Nat} is in the
environment.
\begin{itemize}
\item \lead{Given} a branch guarded only on \texttt{Nat.pred}, \texttt{Nat.add}, \texttt{Nat.mul},
\texttt{Nat.div} or \texttt{Nat.mod};
\item \lead{Then} it still learns \texttt{Nat} is present, from that primitive's own reflection;
dually a guard on \texttt{Nat.bitwise} alone yields \texttt{Bool}, which is how
\texttt{Nat.land}, \texttt{Nat.lor} and \texttt{Nat.xor} obtain it.
\end{itemize}
\end{lemma}
\begin{proof}\leanok
\uses{fam:strong-inversions, thm:isdefeq-foralle-inv, thm:isdefeq-istype, fam:ordered-entrywf}
\lead{Idea} Take \texttt{isType} of the hypothesis, invert the $\Pi$ with
\texttt{IsType.forallE\_inv}, and read the head constant off the domain with
\texttt{HasType.const\_inv}. \texttt{boolOfBitwise} first replaces the recorded type using
\texttt{HasPrimitives.natBitwise}, then peels two $\Pi$s.
\end{proof}

\begin{lemma}[The fixed measure lambda evaluates at ground arguments]
\label{vpc:lem:prim-natfstlamapp}\lean{Lean4Lean.VEnv.HasPrimitives.natFstLamApp}\leanok
\uses{vpc:thm:prim-nat-bool-typing}
\contrib{mixed}\srcloc{Lean4Lean/Verify/Primitive.lean}{96}
\thesisref{The beta rule, axioms.tex, The core language.}

\lead{In short} In any well-formed $\Gamma$, $(\lambda\, m.\ \lambda\, \_.\ m)\ \overline{x}\
\overline{y} \equiv \overline{x}$ -- the \hl{measure} \texttt{unfoldNatWellFounded.WF2} hands the
recogniser for \texttt{Nat.gcd} and \texttt{Nat.bitwise}.
\end{lemma}
\begin{proof}\leanok
\uses{vpc:lem:prim-closedn-natlit, thm:isdefeq-istype}
\lead{Idea} Two \texttt{VEnv.IsDefEq.beta} steps; closedness of the numerals discards the lift and
the final instantiation.
\end{proof}

\begin{lemma}[Nat and Bool translate to their model abbreviations]
\label{vpc:thm:prim-trnat-trbool}
\lean{Lean4Lean.VEnv.HasPrimitives.trNat, Lean4Lean.VEnv.HasPrimitives.trBool}\leanok
\contrib{mixed}\srcloc{Lean4Lean/Verify/Primitive.lean}{122}
\thesisref{Constant typing, axioms.tex, delta reduction.}

\lead{In short} If the environment has the primitives and contains \texttt{Nat} (resp.\
\texttt{Bool}), the source constant \texttt{Nat} (resp.\ \texttt{Bool}) with no universe
arguments translates to $\mathtt{nat}$ (resp.\ $\mathtt{bool}$), at any level and local context.
\end{lemma}
\begin{proof}\leanok
\uses{vpc:thm:prim-nat-bool-typing, fam:strong-inversions}
\lead{Idea} From \texttt{natIsType}/\texttt{boolIsType} at the empty context, read the constant
info off \texttt{HasType.const\_inv} and rebuild with \texttt{TrExprS.const}.
\end{proof}

\begin{lemma}[Translating the checked Nat arrow types]
\label{vpc:thm:prim-natarrow}
\lean{Lean4Lean.TrExprS.natArrow1, Lean4Lean.TrExprS.natArrow2}\leanok
\uses{family:trterm-builders}
\contrib{mixed}\srcloc{Lean4Lean/Verify/Primitive.lean}{149}

\lead{In short} $\forall n_1 : \mathtt{Nat}.\ \mathit{cod}$ and $\forall n_1\,n_2 :
\mathtt{Nat}.\ \mathit{cod}$ -- generic in binder names, \texttt{mdata} and binder info --
translate to $\mathtt{nat} \to \mathit{cod}'$ and $\mathtt{nat} \to \mathtt{nat} \to
\mathit{cod}'$.
\lead{Caveat} Genericity in \texttt{mdata} is what lets these serve the prelude's
strict-implicit declarations verbatim.
\end{lemma}
\begin{proof}\leanok
\uses{vpc:thm:prim-trnat-trbool, vpc:thm:prim-nat-bool-typing}
\lead{Idea} Assemble the arrow from \texttt{TrTy.forallE} and \texttt{TrTy.mdata} over
\texttt{trNat} and \texttt{natIsType} at each binder depth, then \texttt{.trS}.
\end{proof}

\begin{lemma}[Shape inversions for String/Char translations]
\label{vpc:thm:prim-trexprs-shape-inv}
\lean{Lean4Lean.TrExprS.const0_inv, Lean4Lean.TrExprS.appChar_inv, Lean4Lean.TrExprS.appChar_inv', Lean4Lean.TrExprS.listCharStringArrow_inv, Lean4Lean.TrExprS.contains_nat_of_natArrow2}\leanok
\uses{def:vlevel-oflevel}
\contrib{mixed}\srcloc{Lean4Lean/Verify/Primitive.lean}{160}

\lead{In short} Four inversions and one containment fact used by the String/Char branches: a
universe-free constant translates to itself at any context; $c_{\{0\}}\ \mathtt{Char}$ translates
to $c_{\{0\}}\ \mathtt{char}$ (the primed form rebuilds this at any context, both halves closed);
$\mathtt{List\ Char} \to \mathtt{String}$ translates to $\mathtt{listChar} \to \mathtt{string}$;
and a translated binary \texttt{Nat} arrow already witnesses \texttt{Nat}'s presence.
\end{lemma}
\begin{proof}\leanok
\lead{Idea} A \texttt{let}-pattern match on the \texttt{TrExprS} derivation, then
\texttt{simp [VLevel.ofLevel]}; \texttt{appChar\_inv'} additionally weakens both closed
sub-typings with \texttt{HasType.weak0}.
\end{proof}

\begin{lemma}[Nat.bitwise's operator is a boolean binary operator]
\label{vpc:thm:prim-bitwise-operand}\lean{Lean4Lean.TrExprS.bitwiseOperand}\leanok
\uses{family:hasprimitives-accessors}
\srcloc{Lean4Lean/Verify/Primitive.lean}{210}
\thesisref{Unique typing, Pi-injectivity; unique.tex thm:utype, thm:1dinv.}

\lead{In short} If $\mathtt{Nat.bitwise}\ \mathit{op}$ translates to $V$, then $V =
\mathtt{Nat.bitwise}\ \mathit{op}'$ for a translation $\mathit{op}'$, and $\vdash \mathit{op}' :
\mathtt{boolOp2}$ -- read off the recorded type, since nothing in the recogniser checks the
operator's own type.
\lead{Caveat} \texttt{IsDefEqU.forallE\_inv} is \texttt{sorry}-backed, so this branch of the
verification was never unconditional.
\end{lemma}
\begin{proof}\leanok
\uses{thm:foralle-inv-derived, fam:typing-meta}
\lead{Idea} Match the application; replace the recorded type by the spec's using
\texttt{HasPrimitives.natBitwise}; \texttt{HasType.uniqU} and \texttt{IsDefEqU.forallE\_inv} take
the domains apart.
\end{proof}

\begin{lemma}[Application congruence for IsDefEqU]
\label{vpc:lem:prim-isdefequ-app-congr}
\lean{Lean4Lean.VEnv.IsDefEqU.appDF', Lean4Lean.VEnv.IsDefEqU.app_arg', Lean4Lean.VEnv.IsDefEqU.app_fun'}\leanok
\srcloc{Lean4Lean/Verify/Primitive.lean}{244}
\thesisref{Congruence, axioms.tex, core language.}

\lead{In short} Congruence of $\equiv$ under application for the type-hiding
\texttt{IsDefEqU}, with the two typings supplied explicitly, plus the two one-sided
specialisations where only the argument, or only the function, moves.
\lead{Caveat} \texttt{IsDefEqU.of\_l} goes through the \texttt{sorry}-backed
\texttt{IsDefEq.uniq}; every node below that chains two definitional equalities inherits it.
\end{lemma}
\begin{proof}\leanok
\uses{fam:defeq-mixing}
\lead{Idea} \texttt{IsDefEqU.of\_l} recovers the typed judgement on each side, then
\texttt{IsDefEq.appDF}.
\end{proof}

\subsection{The recurrence schemas}

Each schema is stated about an opaque \texttt{VExpr} $f$; the constant enters only afterward,
through \texttt{congr\_head}. All five take \texttt{henv : env.WF}, \texttt{env.HasPrimitives},
containment of \texttt{Nat} or \texttt{Bool}, and \texttt{hU : U = 0} -- the probe context's
level count is only propositionally zero, and this hypothesis lets a caller hand over its facts
unrewritten.

\begin{theorem}[Reflection: unary-step recurrence]
\label{vpc:thm:reflects-unary-step}
\lean{Lean4Lean.VEnv.ReflectsNatNatNat'.of_unary_step_equations}\leanok
\contrib{mixed}\srcloc{Lean4Lean/Verify/Primitive.lean}{273}

\lead{In short} Given $\vdash f : \mathtt{nat} \to \mathtt{nat} \to \mathtt{nat}$ and a closed
unary $g$ evaluating as $g\ \overline{k} \equiv \overline{G\,k}$:
\begin{itemize}
\item \lead{Given} the checker verified $f\ (\mathtt{bvar}\,0)\ \mathtt{zero} \equiv
\mathtt{bvar}\,0$ and $f\ (\mathtt{bvar}\,1)\ (\mathtt{succ}\ (\mathtt{bvar}\,0)) \equiv g\ (f\
(\mathtt{bvar}\,1)\ (\mathtt{bvar}\,0))$, and $F\,a\,0 = a$, $F\,a\,(b+1) = G(F\,a\,b)$;
\item \lead{Then} $f$ \hl{reflects} $F$.
\end{itemize}
\axfoot{sorryAx via VEnv.WF.orderedStrong / VEnv.WF.patsStrong, and via IsDefEq.uniq (UniqueTyping/Injectivity)}
\end{theorem}
\begin{proof}\leanok
\uses{thm:isdefeq-instn, vpc:lem:prim-isdefequ-app-congr, vpc:thm:prim-nat-bool-typing, fam:defeq-mixing}
\lead{Idea} Induction on $b$: the base case instantiates the zero equation with
\texttt{IsDefEqU.instN}; the step instantiates the successor equation twice, closedness of $f$
and $g$ discarding the instantiations, and chains the induction hypothesis under $g$.
\end{proof}

\begin{theorem}[Reflection: binary-step recurrence]
\label{vpc:thm:reflects-binary-step}
\lean{Lean4Lean.VEnv.ReflectsNatNatNat'.of_binary_step_equations}\leanok
\contrib{mixed}\srcloc{Lean4Lean/Verify/Primitive.lean}{312}

\lead{In short} As above, but the successor equation is $f\ x\ (\mathtt{succ}\ y) \equiv g\ (f\ x\
y)\ e$ with $g$ reflecting a binary $G$, and the zero equation returns an open term $z$.
\begin{itemize}
\item \lead{Given} $z[\overline{a}] = \overline{Z(a)}$, $e[\overline{b}][\overline{a}] =
\overline{E(a)}$, $F\,a\,0 = Z(a)$, $F\,a\,(b+1) = G\,(F\,a\,b)\,(E\,a)$;
\item \lead{Then} $f$ reflects $F$.
\end{itemize}
\axfoot{sorryAx via VEnv.WF.orderedStrong / VEnv.WF.patsStrong, and via IsDefEq.uniq (UniqueTyping/Injectivity)}
\end{theorem}
\begin{proof}\leanok
\uses{thm:isdefeq-instn, vpc:lem:prim-isdefequ-app-congr, vpc:thm:prim-nat-bool-typing}
\lead{Idea} Induction on $b$, two \texttt{instN} instantiations per step; the induction
hypothesis moves under $g$ by \texttt{app\_fun'} composed with \texttt{app\_arg'}.
\end{proof}

\begin{lemma}[Reflection: first-argument-step recurrence]
\label{vpc:thm:reflects-first-arg-step}
\lean{Lean4Lean.VEnv.ReflectsNatNatNat'.of_first_arg_step_equations}\leanok
\contrib{mixed}\srcloc{Lean4Lean/Verify/Primitive.lean}{350}

\lead{In short} If the successor equation recurses on a modified first argument, $f\ x\
(\mathtt{succ}\ y) \equiv f\ (g\ \overline{C}\ x)\ y$ with $g$ reflecting $G$, and $F\,a\,0 = a$,
$F\,a\,(b+1) = F\,(G\,C\,a)\,b$, then $f$ reflects $F$.
\axfoot{sorryAx via VEnv.WF.orderedStrong / VEnv.WF.patsStrong, and via IsDefEq.uniq (UniqueTyping/Injectivity)}
\end{lemma}
\begin{proof}\leanok
\uses{thm:isdefeq-instn, vpc:lem:prim-isdefequ-app-congr, vpc:thm:prim-nat-bool-typing}
\lead{Idea} Induction on $b$ generalising the first argument, otherwise the same two-step
\texttt{instN} pattern.
\end{proof}

\begin{lemma}[Reflection: four constructor cases]
\label{vpc:thm:reflects-ctor-cases}
\lean{Lean4Lean.VEnv.ReflectsNatNatBool'.of_constructor_cases}\leanok
\contrib{mixed}\srcloc{Lean4Lean/Verify/Primitive.lean}{386}

\lead{In short} Given $f$'s value at $(0,0)$, $(0,\mathtt{succ}\,y)$, $(\mathtt{succ}\,x,0)$ as
boolean literals and a diagonal $f\ (\mathtt{succ}\,x)\ (\mathtt{succ}\,y) \equiv f\ x\ y$,
matching the four equations for $F$, then $f$ reflects $F$.
\axfoot{sorryAx via VEnv.WF.orderedStrong / VEnv.WF.patsStrong, and via IsDefEq.uniq (UniqueTyping/Injectivity)}
\end{lemma}
\begin{proof}\leanok
\uses{thm:isdefeq-instn, vpc:thm:prim-nat-bool-typing, fam:defeq-mixing}
\lead{Idea} Induction on $a$ generalising $b$, then a four-way case split on $b$: the $(0,0)$ case
is the probed equation verbatim, the mixed cases one \texttt{instN} each, the diagonal case two
instantiations then the induction hypothesis.
\end{proof}

\begin{lemma}[Reflection: boolean binary operator]
\label{vpc:thm:reflects-bool-cases}
\lean{Lean4Lean.VEnv.ReflectsBoolBoolBool'.of_left_cases, Lean4Lean.VEnv.ReflectsBoolBoolBool'.of_closed_cases}\leanok
\contrib{mixed}\srcloc{Lean4Lean/Verify/Primitive.lean}{429}

\lead{In short} Two ways a \texttt{boolOp2} term acquires its reflection: from two open equations
at \texttt{false}/\texttt{true} in the first argument whose right-hand sides are literals
(\texttt{Nat.land}, \texttt{Nat.lor}), or from all four closed combinations at once
(\texttt{Nat.xor}).
\axfoot{sorryAx via VEnv.WF.orderedStrong / VEnv.WF.patsStrong, and via IsDefEq.uniq (UniqueTyping/Injectivity)}
\end{lemma}
\begin{proof}\leanok
\uses{thm:isdefeq-instn, vpc:thm:prim-nat-bool-typing}
\lead{Idea} \texttt{of\_left\_cases} instantiates each equation at the boolean literal,
discarding the instantiation on $f$ by closedness; \texttt{of\_closed\_cases} repackages its
hypotheses after \texttt{subst}.
\end{proof}

\begin{lemma}[Reflection of Nat.pred]
\label{vpc:thm:reflects-pred-eqns}
\lean{Lean4Lean.VEnv.ReflectsNatNat'.of_pred_equations}\leanok
\contrib{mixed}\srcloc{Lean4Lean/Verify/Primitive.lean}{458}

\lead{In short} If $\vdash f : \mathtt{nat} \to \mathtt{nat}$, $f\ \mathtt{zero} \equiv
\mathtt{zero}$ and $f\ (\mathtt{succ}\ (\mathtt{bvar}\,0)) \equiv \mathtt{bvar}\,0$, then $f$
reflects \texttt{Nat.pred}.
\axfoot{sorryAx via VEnv.WF.orderedStrong / VEnv.WF.patsStrong, and via IsDefEq.uniq (UniqueTyping/Injectivity)}
\end{lemma}
\begin{proof}\leanok
\uses{thm:isdefeq-instn, vpc:thm:prim-nat-bool-typing}
\lead{Idea} Case split on the numeral; the successor case is one \texttt{instN} of the open
equation. No induction is needed, the two probes being already the two cases.
\end{proof}

\subsection{From the value to the constant}

\begin{theorem}[Reflection transfers along a definitional equation]
\label{vpc:thm:reflects-congr-head}
\lean{Lean4Lean.VEnv.ReflectsNatNatNat'.congr_head, Lean4Lean.VEnv.ReflectsNatNatBool'.congr_head, Lean4Lean.VEnv.ReflectsNatNat'.congr_head}\leanok
\srcloc{Lean4Lean/Verify/Primitive.lean}{475}
\thesisref{The delta rule, axioms.tex, delta reduction.}

\lead{In short} If $g \equiv f$, $f$ reflects $F$, and $g$ has the expected arrow type, then $g$
\hl{reflects} $F$ -- the only place a constant enters the reflection story.
\lead{Caveat} The numeral typings are a hypothesis, not read from \texttt{HasPrimitives}:
the caller has them only in the extended environment.
\end{theorem}
\begin{proof}\leanok
\uses{vpc:lem:prim-isdefequ-app-congr, fam:defeq-mixing}
\lead{Idea} Two applications of \texttt{IsDefEqU.app\_fun'} move the head across the equation at
each argument, then \texttt{IsDefEqU.trans} with the value's reflection.
\end{proof}

\begin{lemma}[addDefEq's recorded equation]
\label{vpc:thm:isdefequ-todefeq}\lean{Lean4Lean.VEnv.isDefEqU_toDefEq}\leanok
\uses{def:vconstval}
\srcloc{Lean4Lean/Verify/Primitive.lean}{514}
\thesisref{The delta rule, axioms.tex, delta reduction.}

\lead{In short} For a definition value with $\mathit{uvars} = 0$ whose recorded equation is well
formed in the extension, the extension satisfies $\vdash c \equiv v$, for constant $c$ and value
$v$.
\end{lemma}
\begin{proof}\leanok
\lead{Idea} \texttt{IsDefEq.extra0} on the left-injection membership of the fresh equation, then
\texttt{simp} away the empty level substitution.
\end{proof}

\begin{lemma}[Reflections are monotone]
\label{vpc:lem:reflects-mono}
\lean{Lean4Lean.VEnv.ReflectsNatNatNat'.mono, Lean4Lean.VEnv.ReflectsNatNatBool'.mono, Lean4Lean.VEnv.ReflectsNatNat'.mono}\leanok
\srcloc{Lean4Lean/Verify/Primitive.lean}{521}

\lead{In short} If $\mathit{env} \le \mathit{env}'$ and $e$ reflects $F$ in $\mathit{env}$, then
$e$ reflects $F$ in $\mathit{env}'$.
\end{lemma}
\begin{proof}\leanok
\uses{thm:isdefeq-mono}
\lead{Idea} Componentwise \texttt{HasType.mono} and \texttt{IsDefEqU.mono}.
\end{proof}

\begin{theorem}[Adding a definition preserves well-formedness]
\label{vpc:thm:adddefeq-wf}\lean{Lean4Lean.VDefVal.addDefEq_wf}\leanok
\uses{def:vdecl-wf}
\srcloc{Lean4Lean/Verify/Primitive.lean}{536}
\thesisref{Admissibility of a definition, axioms.tex, delta reduction.}

\lead{In short} If $\mathit{venv}$ is well formed, the universe-free definition is well formed
there, and adding its constant succeeds, then adding its equation too yields a well-formed
environment with $\vdash c \equiv v$ at the empty context.
\end{theorem}
\begin{proof}\leanok
\uses{vpc:thm:isdefequ-todefeq}
\lead{Idea} Both halves are the \texttt{.def} step of \texttt{VDecl.WF}; the equation is
\texttt{isDefEqU\_toDefEq} at that step, with \texttt{instL\_id} and
\texttt{VEnv.addConst\_self}.
\end{proof}

\begin{theorem}[Carrying a reflection onto the constant]
\label{vpc:thm:reflects-toconst}
\lean{Lean4Lean.VEnv.ReflectsNatNatNat'.toConst, Lean4Lean.VEnv.ReflectsNatNatBool'.toConst, Lean4Lean.VEnv.ReflectsNatNat'.toConst}\leanok
\contrib{mixed}\srcloc{Lean4Lean/Verify/Primitive.lean}{555}

\lead{In short} If the value reflects $F$ in the checking environment, and that environment is
below the well-formed one the declaration is added to, then the \hl{constant} reflects $F$ after
adding the constant and its equation: the point where every primitive's specification is
established.
\axfoot{sorryAx via VEnv.WF.orderedStrong / VEnv.WF.patsStrong, and via IsDefEq.uniq (UniqueTyping/Injectivity)}
\end{theorem}
\begin{proof}\leanok
\uses{vpc:lem:reflects-mono, vpc:thm:adddefeq-wf, vpc:thm:reflects-congr-head, vpc:thm:prim-nat-bool-typing, fam:defeq-mixing}
\lead{Idea} Three steps carry the reflection across the extension and the definitional equation.
\lead{Steps}
\begin{enumerate}
\item \texttt{.mono} transports the reflection along the composite extension;
\item \texttt{VDefVal.addDefEq\_wf} supplies well-formedness and $c \equiv v$;
\item \texttt{congr\_head} moves the reflection across that equation, typing from
\texttt{HasType.defeqU\_l}.
\end{enumerate}
\end{proof}

\section{The checker-side vocabulary}

This section covers \srcloc{Lean4Lean/Verify/Environment/Primitive/Basic.lean}{1}: vocabulary
every branch shares. Nothing in it names a particular primitive.

\begin{lemma}[A reserved name present in the kernel is present in the model]
\label{vpc:thm:contains-primitive}\lean{Lean4Lean.TypeChecker.VContext.contains_primitive}\leanok
\srcloc{Lean4Lean/Verify/Environment/Primitive/Basic.lean}{19}

\lead{In short} If the checker runs at safety \texttt{safe}, the real environment contains a
reserved primitive name $n$, then the model environment contains $n$.
\end{lemma}
\begin{proof}\leanok
\uses{family:tr-env-find}
\lead{Idea} Turn \texttt{Environment.contains} into an existential \texttt{find?} via
\texttt{SMap.find?\_isSome} and \texttt{map\_wf}, then \texttt{TrEnv.find?\_iff}, whose safety
side condition is \texttt{VContext.safePrimitives}.
\end{proof}

\begin{lemma}[A reflected primitive computes at literals in any context]
\label{vpc:lem:natbinlit}
\lean{Lean4Lean.VEnv.ReflectsNatNatNat.applyLit, Lean4Lean.TypeChecker.VContext.natBinLit, Lean4Lean.TypeChecker.VContext.natBinLitBool}\leanok
\uses{ind:primspec}
\srcloc{Lean4Lean/Verify/Environment/Primitive/Basic.lean}{35}
\thesisref{Weakening, typesys.tex thm:weak.}

\lead{In short} If $c$ carries a \texttt{ReflectsNatNatNat} (resp.\ \texttt{ReflectsNatNatBool})
spec and is present, $\Gamma \vdash c\ \overline{a}\ \overline{b} \equiv \overline{f(a,b)}$ at any
context and universe count, not only the empty one \texttt{HasPrimitives} records.
\end{lemma}
\begin{proof}\leanok
\uses{thm:isdefeq-instl, thm:isdefeq-weakn}
\lead{Idea} Instantiate the universe-free equation with \texttt{instL} at the empty level list
and weaken with \texttt{IsDefEq.weak0}.
\end{proof}

\begin{theorem}[Introducing a Nat- or Bool-typed probe variable]
\label{vpc:thm:withprobe}
\lean{Lean4Lean.TypeChecker.M.WF.withNatProbe, Lean4Lean.TypeChecker.M.WF.withBoolProbe}\leanok
\contrib{mixed}\srcloc{Lean4Lean/Verify/Environment/Primitive/Basic.lean}{63}

\lead{In short} A specialisation of \texttt{M.WF.withLocalDecl} to the only two probe types the
recogniser binds, discharging both side conditions once for the eighteen branches and returning
the fresh variable's translation as \texttt{bvar 0}.
\axfoot{sorryAx via VEnv.WF.orderedStrong / VEnv.WF.patsStrong, invoked directly for the two obligations, and independently via TypeChecker.M.WF.withLocalDecl (Verify/TypeChecker/Basic.lean:489), which is sorry-tainted on master already}
\end{theorem}
\begin{proof}\leanok
\uses{thm:vtc-withlocaldecl, vpc:thm:prim-trnat-trbool, vpc:thm:prim-nat-bool-typing, family:trterm-builders}
\lead{Idea} \texttt{M.WF.withLocalDecl} with \texttt{trNat}/\texttt{natIsType} (resp.\
\texttt{trBool}/\texttt{boolIsType}); the variable's translation is
\texttt{VLCtx.find?\_vlam\_self}.
\end{proof}

\section{Spines, telescopes and substitutions}

\begin{definition}[Spine application and lambda telescopes]
\label{vpc:def:vexpr-appn-lams}
\lean{Lean4Lean.VExpr.appN, Lean4Lean.VExpr.appN_nil, Lean4Lean.VExpr.insts_appN, Lean4Lean.VExpr.subst_appN, Lean4Lean.VExpr.appN_append, Lean4Lean.VExpr.lams, Lean4Lean.VExpr.lams_nil, Lean4Lean.VExpr.lams_cons, Lean4Lean.VExpr.liftN_lams', Lean4Lean.VExpr.liftN_lams, Lean4Lean.VExpr.lams_append}\leanok
\uses{def:subst, def:insts, def:liftn}
\srcloc{Lean4Lean/Verify/Environment/Primitive/Basic.lean}{91}
\thesisref{Telescope notation, axioms.tex sec:inductive.}

\lead{In short} $e.\mathtt{appN}\ [a_1,\dots,a_n] = e\,a_1\cdots a_n$ (leftmost first) and
$\mathtt{lams}\ [A_1,\dots,A_n]\ e = \lambda A_1.\cdots\lambda A_n.\ e$ (outermost first), with
their commutation with substitution, instantiation, lifting and append. \lead{Caveat}
\texttt{liftN\_lams} gives only an existential for the lifted domains; \texttt{liftN\_lams'}
computes them.
\end{definition}

\begin{definition}[Extending a substitution by a value list]
\label{vpc:def:subst-consn}
\lean{Lean4Lean.VExpr.Subst.consN, Lean4Lean.VExpr.Subst.consN_nil, Lean4Lean.VExpr.Subst.consN_cons, Lean4Lean.VExpr.Subst.consN_add, Lean4Lean.VExpr.Subst.consN_append, Lean4Lean.VExpr.Subst.consN_append_singleton, Lean4Lean.VExpr.Subst.tail_cons, Lean4Lean.VExpr.Subst.lift_comp_one, Lean4Lean.VExpr.liftN_subst_consN, Lean4Lean.VExpr.subst_consN_bvars, Lean4Lean.VExpr.subst_lift_tail_inst, Lean4Lean.VExpr.ClosedN.subst_eq'}\leanok
\uses{def:subst, def:liftn}
\srcloc{Lean4Lean/Verify/Environment/Primitive/Basic.lean}{310}
\thesisref{Properties of substitution, typesys.tex thm:subst.}

\lead{In short} $\sigma.\mathtt{consN}\ [v_1,\dots,v_n]$ prepends $n$ values in telescope order
(the last on $\mathtt{bvar}\,0$): a term lifted past $n$ binders ignores them on substitution, the
first $n$ bound variables go to them in telescope order, substituting under a binder then
instantiating composes to a \texttt{cons}, and a closed term is unaffected by any substitution.
\end{definition}

\begin{definition}[Source-level spine application and lambda arity]
\label{vpc:def:expr-spine}
\lean{Lean.Expr.lambdaArity, Lean.Expr.lambdaArity_eq_zero, Lean.Expr.appN, Lean.Expr.appN_eq_mkAppList, Lean.Expr.mkAppN_eq, Lean4Lean.FVarsIn.appN}\leanok
\srcloc{Lean4Lean/Verify/Environment/Primitive/Basic.lean}{214}

\lead{In short} \texttt{Expr.lambdaArity} counts leading lambdas -- how many binders
\texttt{lambdaTelescope} opens, matching \texttt{.lam} syntactically with no \texttt{whnf};
\texttt{Expr.appN} is the list form of \texttt{mkAppN}; \texttt{FVarsIn} distributes over a
spine.
\end{definition}

\begin{lemma}[Inverting spine and free-variable translations]
\label{vpc:thm:trexprs-appn-inv}
\lean{Lean4Lean.TrExprS.appN_inv, Lean4Lean.TrExprS.app2_inv, Lean4Lean.TrExprS.closedN_nil, Lean4Lean.TrExprS.fvar_lift_uniq, Lean4Lean.TrExprS.underBV, Lean4Lean.TrExprS.fvar_uniq}\leanok
\uses{def:vt-vlctx}
\srcloc{Lean4Lean/Verify/Environment/Primitive/Basic.lean}{264}

\lead{In short} A translation of $f\,a_1\cdots a_n$ decomposes into a translation of $f$ and a
\texttt{Forall2} of translated arguments; a translation at the empty context is closed; a free
variable's translation is determined by the context and lifts one binder further;
\texttt{underBV} lifts an arbitrary translation under one more \texttt{vlam} binder.
\end{lemma}
\begin{proof}\leanok
\lead{Idea} Induction over the spine matching \texttt{TrExprS.app}; \texttt{fvar\_uniq} compares
two \texttt{VLCtx.find?} results; \texttt{closedN\_nil} uses well-formedness of the translation;
\texttt{underBV} weakens with a \texttt{BVLift}.
\end{proof}

\begin{definition}[List and Forall2 utilities]
\label{vpc:def:list-helpers}
\lean{List.Forall₂.map_right', List.Forall₂.forall_left, Lean4Lean.List.Forall₂.append_inv, Lean4Lean.List.forall_mem_pair, Lean4Lean.List.Forall₂.snoc, Lean4Lean.List.Forall₂.rev, Lean4Lean.List.Forall₂.fvars_uniq, Lean4Lean.List.mapIdx_replicate_nat}\leanok
\uses{vpc:thm:trexprs-appn-inv}
\srcloc{Lean4Lean/Verify/Environment/Primitive/Basic.lean}{222}

\lead{In short} Routine list-induction lemmas: mapping one side of a \texttt{Forall2}, splitting
over append/snoc/reverse, uniqueness of translated free-variable lists, and that
\texttt{mapIdx} over a replicated \texttt{Nat} telescope is the identity.
\end{definition}

\begin{definition}[A typing context viewed as a VLCtx]
\label{vpc:def:vlctx-ofctx}
\lean{Lean4Lean.VLCtx.ofCtx, Lean4Lean.VLCtx.toCtx_ofCtx}\leanok
\uses{def:vt-vlctx}
\srcloc{Lean4Lean/Verify/Environment/Primitive/Basic.lean}{150}

\lead{In short} Every list of model types is the \texttt{toCtx} of a \texttt{VLCtx} of anonymous
\texttt{vlam} entries, and the round trip is the identity -- transferring \texttt{VLCtx}-only
facts to an arbitrary context.
\end{definition}

\begin{lemma}[Nat and Bool are types, at a plain context]
\label{vpc:thm:prim-istype-prime}
\lean{Lean4Lean.VEnv.HasPrimitives.boolIsType', Lean4Lean.VEnv.HasPrimitives.natIsType'}\leanok
\contrib{mixed}\srcloc{Lean4Lean/Verify/Environment/Primitive/Basic.lean}{158}

\lead{In short} \texttt{natIsType} and \texttt{boolIsType} restated for an ordinary
list-of-types context.
\end{lemma}
\begin{proof}\leanok
\uses{vpc:thm:prim-nat-bool-typing, vpc:def:vlctx-ofctx}
\lead{Idea} Transport along \texttt{VLCtx.toCtx\_ofCtx} and apply the \texttt{VLCtx} version at
\texttt{ofCtx} of the context.
\end{proof}

\begin{lemma}[Both sides of a well-formed application are well formed]
\label{vpc:thm:wf-app-inv2}\lean{Lean4Lean.VExpr.WF.app_inv₂}\leanok
\contrib{mixed}\srcloc{Lean4Lean/Verify/Environment/Primitive/Basic.lean}{431}
\thesisref{Regularity, typesys.tex thm:reg.}

\lead{In short} If $f\,a$ is well formed in $\Gamma$ then so are $f$ and $a$: the type-free form
of \texttt{app\_inv}, which chains as a caller peels a spine apart.
\end{lemma}
\begin{proof}\leanok
\uses{fam:strong-inversions}
\lead{Idea} \texttt{VExpr.WF.app\_inv}, forgetting the two types.
\end{proof}

\begin{lemma}[Spine inversion and spine congruence]
\label{vpc:thm:wf-appn-inv}
\lean{Lean4Lean.VExpr.WF.appN_inv, Lean4Lean.VEnv.IsDefEqU.appN}\leanok
\contrib{mixed}\srcloc{Lean4Lean/Verify/Environment/Primitive/Basic.lean}{437}

\lead{In short} If $e\,a_1\cdots a_n$ is well formed then so is $e$; and if $f \equiv g$ and
$f\,a_1\cdots a_n$ is well formed, then $f\,a_1\cdots a_n \equiv g\,a_1\cdots a_n$ -- the
arguments' typings are recovered, not assumed.
\axfoot{IsDefEqU.appN takes env.WF and derives OrderedStrong from it, so its footprint has both sources: sorryAx via IsDefEqU.of\_l / IsDefEq.uniq (UniqueTyping/Injectivity) and via VEnv.WF.orderedStrong / VEnv.WF.patsStrong. VExpr.WF.appN\_inv takes OrderedStrong as a hypothesis and is itself sorryAx-free}
\end{lemma}
\begin{proof}\leanok
\uses{fam:strong-inversions, fam:defeq-mixing}
\lead{Idea} Induction on the argument list, peeling one application with \texttt{app\_inv}; the
congruence additionally uses \texttt{IsDefEqU.of\_l} and \texttt{IsDefEq.appDF}.
\end{proof}

\begin{lemma}[Inversions onto the head's recorded domain]
\label{vpc:thm:wf-inv-prime}
\lean{Lean4Lean.VExpr.WF.app_inv', Lean4Lean.VExpr.WF.lam_inv', Lean4Lean.VExpr.WF.betaU}\leanok
\contrib{mixed}\srcloc{Lean4Lean/Verify/Environment/Primitive/Basic.lean}{986}

\lead{In short} \texttt{app\_inv'} moves the argument onto the head's \emph{recorded} $\Pi$
domain, usable in a substitution, and types the application; \texttt{lam\_inv'} peels one binder
off a well-formed lambda; \texttt{betaU} gives the argument's typing and the beta equation.
\lead{Caveat} Of the three, only \texttt{lam\_inv'} is \texttt{sorryAx}-free, being the one
that does not go through unique typing.
\axfoot{sorryAx via IsDefEq.uniq and IsDefEqU.forallE\_inv (UniqueTyping/Injectivity)}
\end{lemma}
\begin{proof}\leanok
\uses{fam:strong-inversions, thm:isdefeq-foralle-inv, fam:defeq-mixing}
\lead{Idea} \texttt{app\_inv}, then \texttt{HasType.uniqU} and \texttt{IsDefEqU.forallE\_inv} to
identify the two domains, then \texttt{HasType.defeqU\_r}; \texttt{betaU} closes with
\texttt{VEnv.IsDefEq.beta}.
\end{proof}

\begin{definition}[Arguments closing a lambda telescope]
\label{vpc:ind:argstyped}
\lean{Lean4Lean.VExpr.ArgsTyped, Lean4Lean.VExpr.ArgsTyped.natTelescope, Lean4Lean.OnCtx.of_append, Lean4Lean.OnCtx.append_right}\leanok
\uses{def:hastype-istype, fam:subst-basic}
\srcloc{Lean4Lean/Verify/Environment/Primitive/Basic.lean}{463}
\thesisref{Telescopes, axioms.tex sec:inductive.}

\lead{In short} $\mathsf{ArgsTyped}\ \mathit{env}\ U\ \Gamma_0\ [A_1,\dots,A_n]\ \sigma\
[v_1,\dots,v_n]$ holds when $\Gamma_0 \vdash v_i : A_i[\sigma.\mathtt{consN}\ v_1\cdots
v_{i-1}]$ for each $i$: a relation, not a \texttt{Forall2}, because the $i$-th domain is
substituted by the first $i-1$ arguments. \texttt{Nat}-typed arguments are \texttt{ArgsTyped}
against an all-\texttt{Nat} telescope, \texttt{Nat} being closed; the two \texttt{OnCtx} lemmas
say an appended context restricts, and outer binders leave a well-formed context well formed.
\end{definition}

\begin{theorem}[Telescope arguments define a substitution]
\label{vpc:thm:argstyped-substeq}
\lean{Lean4Lean.VExpr.ArgsTyped.substEq, Lean4Lean.VExpr.wf_lams}\leanok
\uses{vpc:ind:argstyped}
\srcloc{Lean4Lean/Verify/Environment/Primitive/Basic.lean}{495}
\thesisref{Properties of substitution, typesys.tex thm:subst.}

\lead{In short} Typed arguments extend a \texttt{Ctx.SubstEq} from $\Gamma$ to $\Gamma_0$ into one
from $\mathit{As}^{R} \mathbin{+\!\!+} \Gamma$ to $\Gamma_0$; and $\mathtt{lams}\ \mathit{As}\ e$
is well formed in $\Gamma$ whenever $e$ is well formed in the extended, well-formed context.
\end{theorem}
\begin{proof}\leanok
\lead{Idea} Induction on the domain list, \texttt{Ctx.SubstEq.cons} at each step;
\texttt{wf\_lams} iterates \texttt{HasType.lam} the same way.
\end{proof}

\begin{definition}[Pi telescopes and telescope typing]
\label{vpc:def:vexpr-forallEs}
\lean{Lean4Lean.VExpr.forallEs, Lean4Lean.VExpr.forallEs_nil, Lean4Lean.VExpr.forallEs_cons, Lean4Lean.VEnv.HasType.lams, Lean4Lean.VEnv.HasType.appN_forallEs}\leanok
\uses{vpc:def:vexpr-appn-lams, vpc:ind:argstyped}
\srcloc{Lean4Lean/Verify/Environment/Primitive/Basic.lean}{525}
\thesisref{Telescope application, axioms.tex sec:inductive.}

\lead{In short} The $\Pi$ telescope matching \texttt{VExpr.lams}: a lambda telescope over a body
of type $B$ has type $\mathtt{forallEs}\ \mathit{As}\ B$; a $\Pi$-telescope term applied to a
matching \texttt{ArgsTyped} spine has the body's type closed at those arguments, by iterating
\texttt{HasType.app} and rewriting with \texttt{Subst.lift\_comp\_one}.
\end{definition}

\begin{lemma}[Spine congruence with moving arguments]
\label{vpc:thm:isdefequ-appn-prime}
\lean{Lean4Lean.VEnv.IsDefEqU.appN', Lean4Lean.VEnv.IsDefEqU.app_arg}\leanok
\srcloc{Lean4Lean/Verify/Environment/Primitive/Basic.lean}{564}
\thesisref{Congruence, axioms.tex, core language.}

\lead{In short} If $f \equiv f'$, the arguments agree pointwise, and $f\,\bar{x}$ is well formed,
then $f\,\bar{x} \equiv f'\,\bar{y}$; plus the one-step congruence in a single argument.
\end{lemma}
\begin{proof}\leanok
\uses{fam:strong-inversions, vpc:thm:wf-appn-inv}
\lead{Idea} Induction on the argument list, \texttt{appDF} at each step with typings from
\texttt{appN\_inv} and \texttt{app\_inv}.
\end{proof}

\begin{lemma}[Translating an application spine]
\label{vpc:thm:trexprs-appn}\lean{Lean4Lean.TrExprS.appN}\leanok
\contrib{mixed}\srcloc{Lean4Lean/Verify/Environment/Primitive/Basic.lean}{580}

\lead{In short} Given a translation of the head, translations of all arguments, and
well-formedness of the model spine, the source spine translates to it -- the typings
\texttt{TrExprS.app} needs come from well-formedness, not an extra hypothesis.
\end{lemma}
\begin{proof}\leanok
\uses{fam:strong-inversions, vpc:thm:wf-appn-inv}
\lead{Idea} Induction on the argument list, \texttt{VExpr.WF.app\_inv} supplying the two typings
at each \texttt{TrExprS.app}.
\end{proof}

\begin{lemma}[A telescope's domains form a well-formed context]
\label{vpc:thm:lams-ctx}\lean{Lean4Lean.VExpr.lams_ctx}\leanok
\contrib{mixed}\srcloc{Lean4Lean/Verify/Environment/Primitive/Basic.lean}{603}

\lead{In short} If $\Gamma$ is well formed and $\mathtt{lams}\ \mathit{As}\ e$ is well formed in
$\Gamma$, then $\mathit{As}^{R} \mathbin{+\!\!+} \Gamma$ is well formed -- read off the telescope
itself, not carried alongside it.
\end{lemma}
\begin{proof}\leanok
\uses{fam:strong-inversions}
\lead{Idea} Structural recursion peeling one binder with \texttt{VExpr.WF.lam\_inv}.
\end{proof}

\begin{theorem}[Beta for a telescope under a closing substitution]
\label{vpc:thm:lams-appn}
\lean{Lean4Lean.VExpr.lams_appN, Lean4Lean.VExpr.lams_appN'}\leanok
\contrib{mixed}\srcloc{Lean4Lean/Verify/Environment/Primitive/Basic.lean}{618}

\lead{In short} $((\mathtt{lams}\ \mathit{As}\ e)[\sigma])\ v_1\cdots v_n \equiv
e[\sigma.\mathtt{consN}\ \bar{v}]$, the arguments extending $\sigma$ to a closing of the
telescope's context.
The unprimed form derives the argument typings from well-formedness; the primed form takes a
built \texttt{ArgsTyped} and returns well-formedness of the left side instead.
\axfoot{sorryAx via IsDefEq.uniq and IsDefEqU.forallE\_inv (UniqueTyping/Injectivity), and via VEnv.WF.orderedStrong / VEnv.WF.patsStrong, which both forms derive from their own env.WF hypothesis}
\end{theorem}
\begin{proof}\leanok
\uses{fam:strong-inversions, vpc:thm:wf-appn-inv, thm:isdefeq-foralle-inv, vpc:def:subst-consn, thm:isdefeq-substdf, vpc:thm:argstyped-substeq}
\lead{Idea} Induction on the argument list. \lead{Steps}
\begin{enumerate}
\item invert the application (\texttt{appN\_inv}, \texttt{app\_inv}) and the lambda
(\texttt{lam\_inv});
\item reconcile the substituted domain with the recorded one by \texttt{uniqU} and
\texttt{forallE\_inv};
\item take one \texttt{IsDefEq.beta};
\item rewrite the substitution composition with \texttt{Subst.lift\_comp\_one} and
\texttt{subst\_subst}, congruence in the rest by \texttt{IsDefEqU.appN}.
\end{enumerate}
\end{proof}

\section{Ground judgements and the environment extension}

\begin{definition}[Ground judgements and telescope closings]
\label{vpc:def:vcontext-ground}
\lean{Lean4Lean.TypeChecker.VContext.HasType₀, Lean4Lean.TypeChecker.VContext.IsDefEqU₀, Lean4Lean.TypeChecker.VContext.WF₀, Lean4Lean.TypeChecker.VContext.Closing}\leanok
\uses{fam:substeq}
\srcloc{Lean4Lean/Verify/Environment/Primitive/Basic.lean}{684}

\lead{In short} Abbreviations for typing, defeq and well-formedness at the \emph{empty} context of
the checking environment, and for a substitution $\gamma$ closing the recogniser's whole local
context, as a \texttt{Ctx.SubstEq} with itself so the nil case pins it to the identity.
\end{definition}

\begin{definition}[The extension a ground judgement is read in]
\label{vpc:struct:vcontext-ext}
\lean{Lean4Lean.TypeChecker.VContext.Ext, Lean4Lean.TypeChecker.VContext.self, Lean4Lean.TypeChecker.VContext.Ext.cast, Lean4Lean.TypeChecker.VContext.Ext.HasType₀, Lean4Lean.TypeChecker.VContext.Ext.IsDefEqU₀, Lean4Lean.TypeChecker.VContext.Ext.WF₀, Lean4Lean.TypeChecker.VContext.Ext.Closing, Lean4Lean.TypeChecker.VContext.Ext.mono, Lean4Lean.TypeChecker.VContext.Ext.monoT, Lean4Lean.TypeChecker.VContext.Ext.monoW, Lean4Lean.TypeChecker.VContext.Ext.monoIsType, Lean4Lean.TypeChecker.VContext.Ext.monoCtx}\leanok
\uses{def:venv-wf, struct:venv-le}
\srcloc{Lean4Lean/Verify/Environment/Primitive/Basic.lean}{706}
\thesisref{Monotonicity of the judgements under extension, taken for granted by the thesis.}

\lead{In short} A \texttt{VContext.Ext} is a well-formed environment above the checking one.
\begin{itemize}
\item every branch but \texttt{Nat.bitwise} takes the reducible identity extension
\texttt{self}; \texttt{Nat.bitwise} cannot, its spec quantifying over every operator a later
environment might supply; the \texttt{mono} family transports judgements along the extension;
\texttt{cast} moves an extension along a change of context that keeps the environment fixed;
\item \lead{Caveat} \texttt{Ext.cast}'s environment equation defaults to \texttt{rfl}.
\end{itemize}
\end{definition}

\begin{theorem}[A reflected primitive at two literal arguments, under a closing]
\label{vpc:thm:ext-natbinlittr}\lean{Lean4Lean.TypeChecker.VContext.Ext.natBinLitTr}\leanok
\contrib{mixed}\srcloc{Lean4Lean/Verify/Environment/Primitive/Basic.lean}{759}

\lead{In short} For a constant $C$ with a \texttt{ReflectsNatNatNat} spec: if $C\ A\ B$
translates to $r$, $r[\sigma]$ is well formed at the extension's empty context, and $A[\sigma]$,
$B[\sigma]$ are worth $\overline{a}$, $\overline{b}$ for \emph{any} well-formed translations of
$A$, $B$, then $r[\sigma]$ is worth $\overline{F(a,b)}$. Here $\sigma$ is an arbitrary
\texttt{VExpr.Subst}, not assumed a \texttt{Closing}, and containment of $C$ comes from the
translation itself, so a recurrence may name \texttt{Nat.add}, \texttt{Nat.div} or
\texttt{Nat.mod} without its branch guarding for them.
\axfoot{sorryAx via IsDefEq.uniq (UniqueTyping/Injectivity) and, through app\_inv₂, via VEnv.WF.patsStrong}
\end{theorem}
\begin{proof}\leanok
\uses{vpc:thm:trexprs-appn-inv, vpc:thm:prim-trexprs-shape-inv, vpc:lem:natbinlit, vpc:thm:isdefequ-appn-prime, vpc:thm:wf-app-inv2}
\lead{Idea} \texttt{app2\_inv} splits the translation, \texttt{const0\_inv} pins the head;
\texttt{applyLit} evaluates the constant at the two literals, and \texttt{IsDefEqU.appN'} moves
the arguments, well-formedness peeled off by Lemma~\ref{vpc:thm:wf-app-inv2}.
\end{proof}

\section{Opening a lambda telescope}

\begin{lemma}[Counting arguments about MLCtx telescopes]
\label{vpc:thm:mlctx-telescope-counting}
\lean{Lean4Lean.MLCtx.dropN_toCtx_length, Lean4Lean.MLCtx.head_vlam, Lean4Lean.MLCtx.mkLambda'_eq_lams}\leanok
\uses{vpc:def:vexpr-appn-lams}
\srcloc{Lean4Lean/Verify/Environment/Primitive/Basic.lean}{788}
\thesisref{Let binders and zeta, axioms.tex, Let binders.}

\lead{In short} Dropping $n$ binders removes at most $n$ context entries (a \texttt{vlet}
contributing none); a telescope whose binders account for exactly the new entries must open with
a \texttt{vlam}, so \texttt{MLCtx.mkLambda'} agrees with \texttt{VExpr.lams} in that case.
\end{lemma}
\begin{proof}\leanok
\lead{Idea} Induction on $n$ with \texttt{omega} on the length arithmetic; the \texttt{vlet} case
yields a contradiction.
\end{proof}

\begin{definition}[The telescope invariant, translation side included]
\label{vpc:struct:lambdatelescope-inv}\lean{Lean4Lean.lambdaTelescope.Inv}\leanok
\uses{ind:vt-fvlift, ind:trexprs}
\srcloc{Lean4Lean/Verify/Environment/Primitive/Basic.lean}{842}

\lead{In short} What a caller opening $n$ binders gets: the domain list has length $n$; the new
context is $\mathit{As}^{R}$ on the old; a \texttt{VLCtx.FVLift} over the $n$ binders; the opened
variables translate to $\mathtt{bvar}\,0,\dots,\mathtt{bvar}\,(n-1)$ innermost first; and the
original term is \texttt{MLCtx.mkLambda'} of the body, which agrees with \texttt{VExpr.lams} only
when every binder is a \texttt{vlam} -- true here, not assumed elsewhere.
\end{definition}

\begin{definition}[The syntactic Nat-binder guard]
\label{vpc:def:natbindertypes}
\lean{Lean.Expr.natBinderTypes, Lean.Expr.natBinderTypes_of_abstract1, Lean4Lean.MLCtx.mkLambda_natBinderTypes}\leanok
\uses{vpc:thm:mlctx-telescope-counting}
\srcloc{Lean4Lean/Verify/Environment/Primitive/Basic.lean}{856}

\lead{In short} \texttt{Expr.natBinderTypes} decides syntactically that every leading lambda
binder type is the constant \texttt{Nat}; abstraction cannot create such a binder type where
there was none; if the term a telescope built passes the test, its opened context entries are
$\mathtt{nat}$ on the nose.
\lead{Caveat} A hypothesis of the well-founded-recursion theorems below, not a check the
recogniser performs; true by \texttt{rfl} at today's call sites but unverified in general.
\end{definition}

\begin{theorem}[Correctness of lambdaTelescope]
\label{vpc:thm:lambdatelescope-wf}
\lean{Lean4Lean.lambdaTelescope.loop.WF, Lean4Lean.lambdaTelescope.WF}\leanok
\uses{vpc:struct:lambdatelescope-inv}
\srcloc{Lean4Lean/Verify/Environment/Primitive/Basic.lean}{899}

\lead{In short} \texttt{lambdaTelescope e k} opens every leading lambda of $e$ as a fresh probe
variable and hands the instantiated body to $k$, whose continuation may assume the
\texttt{lambdaTelescope.Inv} bundle.
\end{theorem}
\begin{proof}\leanok
\uses{thm:vtc-withlocaldecl, family:trexprs-uninstantiate, family:vt-fvarsin-lemmas, family:expr-subst, thm:trexprs-weakfv}
\lead{Idea} Induction via \texttt{loop.WF}: each iteration matches the \texttt{lam} rule, applies
\texttt{M.WF.withLocalDecl}, shows the new abstraction cancels the instantiation, rebuilds the
\texttt{Inv} fields, and opens the binder with \texttt{TrExprS.inst\_fvar}.
\end{proof}

\section{What a branch is given, and what it owes}

\begin{definition}[What a branch is given about the definition]
\label{vpc:struct:primitive-data}
\lean{Lean4Lean.Primitive.Data, Lean4Lean.Primitive.Data.contains, Lean4Lean.Primitive.Data.uvars_eq, Lean4Lean.Primitive.Data.hv}\leanok
\uses{ind:trexprs, def:vconstval, vpc:thm:contains-primitive, family:trterm-builders}
\srcloc{Lean4Lean/Verify/Environment/Primitive/Basic.lean}{1036}
\thesisref{Admissibility of a definition, axioms.tex, delta reduction.}

\lead{In short} A record packaging what \texttt{addDefinition} has already established before
calling the recogniser: safety, the level parameters, matching arities, and typed translations of
the definition's type and value; and three derived facts: a named primitive present in the real
environment is in the model, the universe count is $0$ once the entry guard fires, and the value
is a bundled \texttt{TrTerm} at any probe context.
\end{definition}

\begin{definition}[The translation atoms shared by all branches]
\label{vpc:def:prim-atoms}
\lean{Lean4Lean.Primitive.zerob, Lean4Lean.Primitive.succb, Lean4Lean.Primitive.oneb, Lean4Lean.Primitive.twob, Lean4Lean.Primitive.boolb, Lean4Lean.Primitive.hNatT, Lean4Lean.Primitive.trNat, Lean4Lean.Primitive.predb, Lean4Lean.Primitive.addb, Lean4Lean.Primitive.divb, Lean4Lean.Primitive.modb, Lean4Lean.Primitive.mulb, Lean4Lean.Primitive.natCod, Lean4Lean.Primitive.natCod1, Lean4Lean.Primitive.boolCod, Lean4Lean.Primitive.trBool, Lean4Lean.Primitive.boolOp2Ty, Lean4Lean.Primitive.boolOp2Ty_tgt, Lean4Lean.Primitive.bitwiseTy}\leanok
\uses{vpc:thm:prim-trnat-trbool, vpc:thm:prim-nat-bool-typing, struct:trterm, thm:trexprs-ofconst, family:trterm-builders}
\contrib{mixed}\srcloc{Lean4Lean/Verify/Environment/Primitive/Basic.lean}{1053}

\lead{In short} Pre-built \texttt{TrTerm}/\texttt{TrTy} bundles for everything a branch's shape
check mentions: the \texttt{Nat} constructors and the numerals $1,2$ as constructor applications,
the boolean literals, the operators named in someone else's recurrence, the codomain bundles at
each depth, \texttt{Bool -> Bool -> Bool} pinned to target \texttt{boolOp2}, and
\texttt{Nat.bitwise}'s own type.
\axfoot{sorryAx via VEnv.WF.orderedStrong / VEnv.WF.patsStrong, entering through trNat, trBool and TrTerm.ofConst}
\end{definition}

\begin{theorem}[The checked type pins the constant's model type]
\label{vpc:thm:data-mktyeq}
\lean{Lean4Lean.Primitive.Data.mkTyEq, Lean4Lean.Primitive.Data.mkTyEq1, Lean4Lean.Primitive.Data.mkTyEqBitwise}\leanok
\contrib{mixed}\srcloc{Lean4Lean/Verify/Environment/Primitive/Basic.lean}{1163}

\lead{In short} If the syntactic test showed the source type has the expected arrow shape, the
model type recorded for the constant is exactly the matching model arrow: $\mathtt{nat} \to
\mathtt{nat} \to \mathit{cod}'$, $\mathtt{nat} \to \mathit{cod}'$, or $\mathtt{boolOp2} \to
\mathtt{natOp2}$, the binary form generic in the codomain bundle so \texttt{Nat.beq}/\texttt{Nat.ble}
pass a \texttt{Bool} bundle instead.
\axfoot{sorryAx via VEnv.WF.orderedStrong / VEnv.WF.patsStrong, inherited from natArrow1/natArrow2 and the atoms}
\end{theorem}
\begin{proof}\leanok
\uses{family:trexprs-eqv, def:trexprs-isunique, vpc:thm:prim-natarrow, vpc:def:prim-atoms}
\lead{Idea} \texttt{TrExprS.eqv} moves the translation onto the compared shape, then
\texttt{TrExprS.unique} against the arrow \texttt{natArrow1}/\texttt{natArrow2}/\texttt{bitwiseTy}
build.
\end{proof}

\begin{definition}[The obligation the recognizer must discharge]
\label{vpc:struct:primitiveresult}\lean{Lean4Lean.Primitive.PrimitiveResult}\leanok
\uses{ind:primspec, def:tr-constant, def:venv-basic}
\srcloc{Lean4Lean/Verify/Environment/Primitive/Basic.lean}{1204}
\thesisref{Conservativity of definitions, axioms.tex, delta reduction.}

\lead{In short} For a definition the recogniser accepted: it is safe, has no level parameters, and
for every well-formed extension already satisfying \texttt{HasPrimitives}, adding its constant
and its equation keeps that invariant.
\end{definition}

\begin{theorem}[Assembling a PrimitiveResult]
\label{vpc:thm:data-mkresult}
\lean{Lean4Lean.Primitive.Data.mkResult, Lean4Lean.Primitive.Data.mkResult1, Lean4Lean.Primitive.Data.mkResultBool, Lean4Lean.Primitive.Data.mkResultTypeEq}\leanok
\uses{vpc:struct:primitiveresult}
\srcloc{Lean4Lean/Verify/Environment/Primitive/Basic.lean}{1216}

\lead{In short} Given membership in \texttt{primSpecs}, the pinned model type, and a reflection
of the value in the checking environment, produce the \texttt{PrimitiveResult}; the
quantification over extensions is discharged here, so no branch sees it.
\texttt{mkResultTypeEq} covers \texttt{Char.ofNat}, whose whole obligation is the recorded
type.
\end{theorem}
\begin{proof}\leanok
\uses{thm:hasprimitives-addconst, vpc:thm:reflects-toconst, vpc:struct:primitive-data}
\lead{Idea} \texttt{HasPrimitives.addPrimitiveDefEq} reduces the invariant to the one new
specification, then the matching \texttt{toConst} carries the value's reflection onto the
constant.
\end{proof}

\begin{theorem}[The obligation for Nat.bitwise]
\label{vpc:thm:data-mkresultbitwise}\lean{Lean4Lean.Primitive.Data.mkResultBitwise}\leanok
\uses{vpc:struct:primitiveresult}
\contrib{mixed}\srcloc{Lean4Lean/Verify/Environment/Primitive/Basic.lean}{1262}

\lead{In short} For \texttt{Nat.bitwise} the reflection must hold at \emph{every} boolean
operator a well-formed extension might supply: given the recorded type and, for every extension
and every $f$ reflecting $g$, that the value applied to $f$ reflects
$\mathtt{Nat.bitwise}\ g$, the \texttt{PrimitiveResult} follows: the specification that forces
\texttt{VContext.Ext} (\ref{vpc:struct:vcontext-ext}) to exist at all.
\axfoot{sorryAx via VEnv.WF.orderedStrong / VEnv.WF.patsStrong (natLitT) and via IsDefEq.uniq (UniqueTyping/Injectivity)}
\end{theorem}
\begin{proof}\leanok
\uses{thm:hasprimitives-addconst, vpc:thm:adddefeq-wf, vpc:thm:reflects-congr-head, vpc:lem:prim-isdefequ-app-congr, vpc:thm:prim-nat-bool-typing, fam:defeq-mixing}
\lead{Idea} \texttt{addPrimitiveDefEq} reduces to the new specification; transport the argument
along the composite extension, then move it from value to constant with \texttt{congr\_head}.
\end{proof}

\section{The Nat.div and Nat.mod fuel recursion}

\begin{definition}[The shape of a fuel-recursion call]
\label{vpc:def:natgocall}
\lean{Lean4Lean.Primitive.natGoCall, Lean4Lean.Primitive.natGoType}\leanok
\srcloc{Lean4Lean/Verify/Environment/Primitive/Basic.lean}{1313}

\lead{In short} A call of \texttt{Nat.div.go} or \texttt{Nat.modCore.go} as a five-argument
application -- divisor, positivity proof, fuel, dividend, a proof the fuel covers it -- and the
type it is checked against.
\end{definition}

\begin{lemma}[Typing the arguments of a fuel-recursion call]
\label{vpc:thm:goargs}\lean{Lean4Lean.Primitive.goArgs}\leanok
\uses{vpc:def:natgocall}
\contrib{mixed}\srcloc{Lean4Lean/Verify/Environment/Primitive/Basic.lean}{1325}

\lead{In short} If a well-formed five-argument call at the empty context has the recorded
\texttt{natGoType} head, each argument has its recorded domain and the call itself has type
$\mathtt{nat}$.
\axfoot{sorryAx via IsDefEq.uniq and IsDefEqU.forallE\_inv through app\_inv', and via VEnv.WF.orderedStrong / VEnv.WF.patsStrong, which the four app\_inv calls invoke directly}
\end{lemma}
\begin{proof}\leanok
\uses{fam:strong-inversions, vpc:thm:wf-inv-prime, fam:inst-inst}
\lead{Idea} Four \texttt{app\_inv}s peel the spine, then five \texttt{app\_inv'}s move each
argument onto its recorded domain, \texttt{simp} discharging each instantiation via closedness of
the order relation.
\end{proof}

\begin{theorem}[The fuel recursion runs to the end]
\label{vpc:thm:natfuelrec}\lean{Lean4Lean.Primitive.natFuelRec}\leanok
\uses{vpc:def:natgocall}
\srcloc{Lean4Lean/Verify/Environment/Primitive/Basic.lean}{1360}

\lead{In short} Given a step equation at literals -- a call at fuel $f'+1$ stops with
$\overline{E(x)}$ when $x < b$, else reduces to a context $H$ applied to the call at fuel $f'$ and
dividend $x - b$ -- with $F,E,h$ arithmetically consistent, a well-formed call at positive $b$ and
fuel $f > x$ is definitionally $\overline{F(x,b)}$.
Positivity of $b$ is a hypothesis, not derived.
\end{theorem}
\begin{proof}\leanok
\uses{vpc:thm:goargs, fam:defeq-mixing}
\lead{Idea} Induction on the fuel, generalising the dividend and the bound proof, chaining the
induction hypothesis through $H$'s evaluation and the argument typings from \texttt{goArgs}.
\end{proof}

\begin{lemma}[Guard and safety boilerplate]
\label{vpc:lem:elsefail-hok}
\lean{Lean4Lean.Primitive.elseFail, Lean4Lean.Primitive.hok}\leanok
\srcloc{Lean4Lean/Verify/Environment/Primitive/Basic.lean}{1392}

\lead{In short} A guard that throws on failure proves whatever its continuation proves; the
recogniser's entry test implies the definition is safe with no level parameters.
\end{lemma}
\begin{proof}\leanok
\uses{thm:vtc-forin}
\lead{Idea} \texttt{split} and \texttt{M.WF.bindThrow}; \texttt{simp\_all} on the entry test.
\end{proof}

\section{Well-founded recursion}

This section covers \srcloc{Lean4Lean/Verify/Environment/Primitive/Recursion.lean}{1}: the
verification of \texttt{unfoldNatWellFounded}, whose only callers guard \texttt{Nat.gcd} and
\texttt{Nat.bitwise} (\texttt{Nat.div}/\texttt{Nat.mod} use the fuel recursion above).
$\mathtt{Nat.fix}\ h\ F\ x$ does not unfold definitionally to $F\ x\ (\lambda y\,\_.\
\mathtt{Nat.fix}\ h\ F\ y)$, so a caller must consume fuel to reach it.

\begin{definition}[What the recognizer hands back]
\label{vpc:struct:probebundle}\lean{Lean4Lean.Primitive.ProbeBundle}\leanok
\uses{family:kernel-primitive-dsl, ind:trexprs, vpc:def:vexpr-appn-lams}
\srcloc{Lean4Lean/Verify/Environment/Primitive/Recursion.lean}{43}
\thesisref{No counterpart; the ambient machinery is axioms.tex sec:inductive.}

\lead{In short} A record extending the code-level \texttt{Probe} with translations of the
fixpoint functional $F$, the packer, and the induction-hypothesis domain; the packed-argument
type $\mathit{Aty}$ read off $F$'s inferred type; that the packer \emph{is} a lambda telescope
over \texttt{Nat} domains with body of type $\mathit{Aty}$; and the typings of $\mathit{dom}$ and
$F$.
\lead{Caveat} Nothing here relates $\mathit{Aty}$ to the packer; that check is pushed to
\texttt{applied}.
\end{definition}

\begin{lemma}[The bundle applied to a packed argument]
\label{vpc:thm:probebundle-applied}
\lean{Lean4Lean.Primitive.ProbeBundle.applied, Lean4Lean.Primitive.ProbeBundle.packAty}\leanok
\uses{vpc:struct:probebundle}
\srcloc{Lean4Lean/Verify/Environment/Primitive/Recursion.lean}{80}

\lead{In short} \texttt{applied}: given a translated argument of type $\mathit{Aty}$, get
translations of $\mathit{dom}\ a$ and $F\ a$, that $\mathit{dom}\ a$ is a type, and that $F\ a$ is
well typed, at the probe's own extended context.
\texttt{packAty}: the packer applied to a \texttt{Nat}-typed spine of the right length has
type $\mathit{Aty}$.
\end{lemma}
\begin{proof}\leanok
\uses{vpc:def:vexpr-forallEs, vpc:ind:argstyped, thm:trexprs-weakfv, vpc:thm:lams-appn}
\lead{Idea} Weaken the bundle's typings over the probe's lift and apply \texttt{HasType.app};
\texttt{packAty} rewrites the packer as a \texttt{Nat} telescope, builds \texttt{ArgsTyped} by
\texttt{natTelescope}, and closes with \texttt{lams\_appN'}.
\end{proof}

\begin{definition}[The bundle under a closing substitution]
\label{vpc:def:probebundle-closed}
\lean{Lean4Lean.Primitive.ProbeBundle.Fc, Lean4Lean.Primitive.ProbeBundle.packc, Lean4Lean.Primitive.ProbeBundle.domc, Lean4Lean.Primitive.ProbeBundle.arg, Lean4Lean.Primitive.ProbeBundle.ihTy, Lean4Lean.Primitive.ProbeBundle.pγ, Lean4Lean.Primitive.ProbeBundle.parg, Lean4Lean.Primitive.ProbeBundle.pihTy, Lean4Lean.Primitive.ProbeBundle.plhs, Lean4Lean.Primitive.ProbeBundle.ihCtx}\leanok
\uses{def:subst, vpc:def:vexpr-appn-lams}
\srcloc{Lean4Lean/Verify/Environment/Primitive/Recursion.lean}{154}

\lead{In short} Accessors for the functional, packer and domain closed by a substitution
$\gamma$; the packed argument and induction-hypothesis type at ground arguments; and the
\texttt{p}-prefixed family giving the same inside a probe, together with its context and the
closed equation's left side. The recursion arguments are already closed, being literals, so
$\gamma$ does not touch them.
\end{definition}

\begin{definition}[One unfolding with recursive calls discharged]
\label{vpc:def:probebundle-done}\lean{Lean4Lean.Primitive.ProbeBundle.Done}\leanok
\uses{vpc:def:probebundle-closed, vpc:struct:vcontext-ext}
\srcloc{Lean4Lean/Verify/Environment/Primitive/Recursion.lean}{176}

\lead{In short} $\mathsf{Done}$: for every recursion argument $a$ and induction hypothesis
$\mathit{ih}$ that returns $\overline{R(y)}$ on every $y$ of strictly smaller measure, $F\
(\mathit{pack}\ \overline{\sigma(a)})\ \mathit{ih} \equiv \overline{R(a)}$ -- the caller's real
obligation, the smaller-measure restriction being all the fuel induction below can supply.
\end{definition}

\begin{definition}[What the checks buy, at ground arguments]
\label{vpc:struct:natfixunfold}\lean{Lean4Lean.Primitive.ProbeBundle.NatFixUnfold}\leanok
\uses{vpc:def:probebundle-closed}
\srcloc{Lean4Lean/Verify/Environment/Primitive/Recursion.lean}{223}

\lead{In short} Three fields over a family $\mathit{goFn}$: the packed argument is well typed;
\texttt{entry} says the value at ground arguments is $\mathit{go}$ at fuel
$\mathit{measure}(a)+1$; \texttt{step} says one unit of fuel buys one call to $F$ with recursive
calls at the lower fuel, uniformly in a base point $b$; $\mathit{go}$ is a family since closing
the telescope at different arguments may leave different terms, only the functional being
telescope-free.
\end{definition}

\begin{theorem}[The fuel induction]
\label{vpc:thm:goconverges}
\lean{Lean4Lean.Primitive.ProbeBundle.GoConverges, Lean4Lean.Primitive.ProbeBundle.GoConverges.succ, Lean4Lean.Primitive.ProbeBundle.GoConverges.all, Lean4Lean.Primitive.ProbeBundle.NatFixUnfold.reflects}\leanok
\uses{vpc:struct:natfixunfold, vpc:def:probebundle-done}
\srcloc{Lean4Lean/Verify/Environment/Primitive/Recursion.lean}{247}

\lead{In short} $\mathsf{GoConverges}\ t$: $\mathit{go}$ at fuel $t$ computes $R$ on every
argument whose measure $t$ outruns; preserved by one unit of fuel given \texttt{NatFixUnfold} and
\texttt{Done}, hence holds at every fuel, giving the reflection at the fuel actually entered.
\end{theorem}
\begin{proof}\leanok
\uses{fam:defeq-mixing}
\lead{Idea} Induction on the fuel: \texttt{succ} feeds \texttt{step}'s hypothesis to
\texttt{Done}, discharging its recursive-call premise with the outer induction hypothesis.
\end{proof}

\begin{theorem}[Correctness of running one probe against the bundle]
\label{vpc:thm:probebundle-probe-wf}\lean{Lean4Lean.Primitive.ProbeBundle.probe.WF}\leanok
\uses{vpc:def:probebundle-closed}
\contrib{mixed}\srcloc{Lean4Lean/Verify/Environment/Primitive/Recursion.lean}{344}

\lead{In short} Running a probe with a substitution, a failure continuation and a
right-hand-side builder is well formed, and its postcondition has the shape \texttt{Done} asks
about: for every extension, closing and closed induction hypothesis, the probed left side is
definitionally equal to a value the builder produced and described.
\axfoot{sorryAx via VEnv.WF.orderedStrong / VEnv.WF.patsStrong (TrExprS.appN) and via the typechecker monad's own gaps}
\end{theorem}
\begin{proof}\leanok
\uses{vpc:thm:probebundle-applied, vpc:thm:trexprs-appn, thm:vtc-withlocaldecl, vpc:thm:prim-istype-prime}
\lead{Idea} Assemble the packed argument's typing and translation via \texttt{packAty} and
\texttt{TrExprS.appN}; \texttt{applied} gives the binder's type; open it with
\texttt{M.WF.withLocalDecl}, run the builder, close by substituting along the extended closing.
\end{proof}

\begin{definition}[The go step equation, one telescope at a time]
\label{vpc:def:unfold-step}
\lean{Lean4Lean.Primitive.unfoldNatWellFounded.StepIH, Lean4Lean.Primitive.unfoldNatWellFounded.StepBranch, Lean4Lean.Primitive.unfoldNatWellFounded.Step}\leanok
\uses{vpc:struct:lambdatelescope-inv, ind:trexprs}
\srcloc{Lean4Lean/Verify/Environment/Primitive/Recursion.lean}{445}
\thesisref{The iota rule for Nat, axioms.tex sec:inductive.}

\lead{In short} Three nested existentials recording what the recogniser's three
\texttt{lambdaTelescope} passes established about $\mathit{go}\ \dots (\mathtt{succ}\ t)\ x\
\mathit{hfuel} \equiv F\ x\ (\lambda y\,h_y.\ \mathit{go}\ \dots t\ y\ \mathit{pf})$:
\texttt{Step} is \texttt{go}'s own telescope (its body a recursor on the fuel variable, which the
recursor does not mention); \texttt{StepBranch} is the successor branch as a function of the
recursion argument; \texttt{StepIH} is the induction hypothesis as the recursor at the lower
fuel.
\end{definition}

\begin{definition}[The eager gadget at literals]
\label{vpc:def:unfold-eager}\lean{Lean4Lean.Primitive.unfoldNatWellFounded.Eager}\leanok
\uses{ind:trexprs}
\srcloc{Lean4Lean/Verify/Environment/Primitive/Recursion.lean}{501}

\lead{In short} What the check on the \texttt{eager} gadget is worth: a translation of the
gadget such that at every literal $k$, $\mathit{eager}\ \overline{k} \equiv \overline{k}$. Not an
implication on the decision being a literal: that premise would live at the probe context, where
the decision is precisely not one.
\end{definition}

\begin{definition}[What the local-declaration block exports]
\label{vpc:def:unfold-blockq}\lean{Lean4Lean.Primitive.unfoldNatWellFounded.BlockQ}\leanok
\uses{vpc:def:unfold-step, vpc:def:unfold-eager, ind:trexprs, def:trexpr}
\srcloc{Lean4Lean/Verify/Environment/Primitive/Recursion.lean}{527}

\lead{In short} The block's four checks, carried forward as a proposition about the extended
context: the fixpoint's measure agrees with the caller's; the entry equation reduces the fixpoint
at the bound variable to \texttt{go} fuelled by the successor of the measure; the \texttt{Step}
equation at that \texttt{go}; and the \texttt{Eager} gadget.
\end{definition}

\begin{theorem}[The recognizer returns a bundle satisfying the unfolding]
\label{vpc:thm:unfoldnatwf-prime}\lean{Lean4Lean.Primitive.unfoldNatWellFounded.WF'}\leanok
\uses{vpc:struct:natfixunfold, vpc:def:natbindertypes}
\srcloc{Lean4Lean/Verify/Environment/Primitive/Recursion.lean}{587}

\lead{In short} \texttt{unfoldNatWellFounded e meas fail} is well formed and returns a
\texttt{ProbeBundle} whose packer telescope has exactly \texttt{meas.lambdaArity} domains, such
that for every extension and closing there is a \texttt{NatFixUnfold} for it at the value.
Hypotheses: translations of the value and measure, presence of \texttt{Nat}, safety, the
conditional-verification facts moved to callers, the measure's evaluation at ground arguments,
and the \texttt{natBinderTypes} guard on the measure.
\end{theorem}
\begin{proof}\leanok
\uses{vpc:thm:lambdatelescope-wf, vpc:def:unfold-blockq, vpc:thm:lams-appn, vpc:thm:argstyped-substeq, thm:vtc-withlocaldecl}
\lead{Idea} Open the measure's telescope with \texttt{lambdaTelescope.WF} and relate it to the
value's via \texttt{checkType}. \lead{Steps}
\begin{enumerate}
\item track translations through \texttt{whnfCore} and \texttt{unfoldDefinition};
\item destructure the application;
\item turn \texttt{BlockQ}'s four checks into \texttt{NatFixUnfold}'s three fields;
\item close the telescope per base point with \texttt{VExpr.lams\_appN} and
\texttt{ArgsTyped.substEq}.
\end{enumerate}
\end{proof}

\begin{theorem}[The recognizer's contract for its callers]
\label{vpc:thm:unfoldnatwf}
\lean{Lean4Lean.Primitive.unfoldNatWellFounded.WF, Lean4Lean.Primitive.unfoldNatWellFounded.WF₂}\leanok
\uses{vpc:def:probebundle-done}
\srcloc{Lean4Lean/Verify/Environment/Primitive/Recursion.lean}{1723}

\lead{In short} Combining the plumbing with the fuel induction: the recogniser returns a bundle
such that for every extension, closing and intended answer $R$, if the caller proves
\texttt{Done}, the value at ground arguments is $\overline{R(a)}$.
The second form specialises to the two-argument recursion measured by $\lambda m\,\_.\ m$,
used by \texttt{Nat.gcd} and \texttt{Nat.bitwise}.
\end{theorem}
\begin{proof}\leanok
\uses{vpc:thm:unfoldnatwf-prime, vpc:thm:goconverges, vpc:lem:prim-natfstlamapp, vpc:thm:prim-trnat-trbool, vpc:def:prim-atoms}
\lead{Idea} Weaken the plumbing theorem's postcondition and apply
\texttt{NatFixUnfold.reflects}; the specialised form builds the fixed measure's translation and
evaluates it with \texttt{natFstLamApp}.
\end{proof}

\section{The verification boundary}

This section covers \srcloc{Lean4Lean/Verify/Environment/Primitive.lean}{1}: the theorem the
declaration checker consumes, and a pre-staged specification of the recogniser's inductive
clause.

\begin{theorem}[Verification boundary for the recognizer]
\label{vpc:thm:checkdef-wf}\lean{Lean4Lean.Primitive.checkDef.WF}\leanok
\uses{vpc:struct:primitive-data, vpc:struct:primitiveresult}
\srcloc{Lean4Lean/Verify/Environment/Primitive.lean}{33}
\thesisref{No counterpart; the Lean 4 kernel's primitive-acceleration layer.}

\lead{In short} If the model environment is well formed for the real one, universe counts agree,
and the definition's type and value already have their recorded translations, then
\texttt{Primitive.checkDef} either returns \texttt{false} or returns \texttt{true} with a
\texttt{PrimitiveResult}: safe, no universe parameters, addition preserves
\texttt{HasPrimitives} in every well-formed extension.
\end{theorem}
\begin{proof}\leanok
\uses{prim-clauses:thm:checknatadd-wf, prim-divmod:thm:checknatdiv-wf, prim-gcd:thm:checknatgcd-wf, prim-bitwise:thm:checknatbitwise-wf, prim-clauses:thm:checkstringoflist-wf}
\lead{Idea} Bundle the hypotheses into a \texttt{Primitive.Data}, unfold \texttt{checkDef} and
split on its eighteen-way dispatch, each branch discharged by its
Chapter~\ref{chap:primitives-arith} theorem; the default branch returns \texttt{nofun}.
\end{proof}

\begin{definition}[Adding a run of constants]
\label{vpc:ind:addsconsts}\lean{Lean4Lean.Primitive.AddsConsts}\leanok
\uses{def:venv-basic}
\srcloc{Lean4Lean/Verify/Environment/Primitive.lean}{73}
\thesisref{Inductive type admissibility, axioms.tex sec:inductive.}

\lead{In short} An inductive relation: one environment is reached from another by adding each
name of a list at its declared universe-free type, interleaved with arbitrary
\texttt{addDefEq} steps.
\lead{Caveat} It abstracts the only part of a future \texttt{AddInduct} specification the
primitive invariant depends on; nothing yet constructs an \texttt{AddsConsts}.
\end{definition}

\begin{lemma}[Basic properties of a run of additions]
\label{vpc:lem:addsconsts-lemmas}
\lean{Lean4Lean.Primitive.AddsConsts.constants_stable, Lean4Lean.Primitive.AddsConsts.constants_of_mem, Lean4Lean.Primitive.AddsConsts.holds_of_ne}\leanok
\uses{vpc:ind:addsconsts}
\srcloc{Lean4Lean/Verify/Environment/Primitive.lean}{81}

\lead{In short} A constant already present survives the run; each added constant is present at
the end with its added type; a primitive specification whose name is none of the added ones
survives the run.
\end{lemma}
\begin{proof}\leanok
\uses{fam:addconsts, thm:primspec-addconst}
\lead{Idea} Induction on the derivation, using \texttt{VEnv.addConst\_constants\_eq} and the
\texttt{PrimSpec.Holds} transport lemmas.
\end{proof}

\begin{definition}[Constants contributed by Bool and Nat]
\label{vpc:def:bool-nat-consts}
\lean{Lean4Lean.Primitive.boolConsts, Lean4Lean.Primitive.natConsts}\leanok
\srcloc{Lean4Lean/Verify/Environment/Primitive.lean}{114}

\lead{In short} The lists
$[(\mathtt{Bool},\mathtt{Type}),(\mathtt{false},\mathtt{bool}),(\mathtt{true},\mathtt{bool})]$
and $[(\mathtt{Nat},\mathtt{Type}),(\mathtt{Nat.zero},\mathtt{nat}),(\mathtt{Nat.succ},\mathtt{nat}
\to \mathtt{nat})]$.
\end{definition}

\begin{lemma}[Adding the Bool or Nat block preserves the invariant]
\label{vpc:thm:addsconsts-hasprimitives}
\lean{Lean4Lean.Primitive.AddsConsts.hasPrimitives_bool, Lean4Lean.Primitive.AddsConsts.hasPrimitives_nat}\leanok
\uses{vpc:def:bool-nat-consts, ind:primspec}
\srcloc{Lean4Lean/Verify/Environment/Primitive.lean}{121}

\lead{In short} If an environment satisfies \texttt{HasPrimitives} and the next is obtained by
adding exactly \texttt{boolConsts} (resp.\ \texttt{natConsts}), it still satisfies
\texttt{HasPrimitives}.
\end{lemma}
\begin{proof}\leanok
\uses{vpc:lem:addsconsts-lemmas}
\lead{Idea} Case split on whether the specification's name is among the three added; if not,
\texttt{holds\_of\_ne}, else decide by membership and discharge with
\texttt{constants\_of\_mem}.
\end{proof}

\begin{definition}[What the inductive clause is worth]
\label{vpc:def:primitive-inductive-result}\lean{Lean4Lean.Primitive.PrimitiveInductiveResult}\leanok
\uses{vpc:def:bool-nat-consts, vpc:ind:addsconsts}
\srcloc{Lean4Lean/Verify/Environment/Primitive.lean}{162}
\thesisref{Inductive specifications, axioms.tex sec:inductive.}

\lead{In short} A flagged inductive declaration is one of the two primitive inductives in
standard form (safe, no level or ordinary parameters, one type at \texttt{Type}, exactly the
constructors the table names) and adding the constants it contributes preserves
\texttt{HasPrimitives}.
\lead{Caveat} The recogniser's inductive clause waits on a constructive
\texttt{AddInduct} supplying exactly these constants; none exists yet.
\end{definition}

\begin{lemma}[Verification boundary for the inductive clause]
\label{vpc:thm:checkinductive-wf}\lean{Lean4Lean.Primitive.checkInductive.WF}\leanok
\uses{vpc:def:primitive-inductive-result}
\srcloc{Lean4Lean/Verify/Environment/Primitive.lean}{174}

\lead{In short} If \texttt{Primitive.checkInductive} returns \texttt{true} on a declaration, then
\texttt{PrimitiveInductiveResult} holds of it.
\lead{Caveat} Unused: the environment checker calls \texttt{checkInductive}, but no
verification of the declaration checker's inductive case consumes this theorem yet.
\end{lemma}
\begin{proof}\leanok
\uses{vpc:thm:addsconsts-hasprimitives, family:expr-beq}
\lead{Idea} Unfold and split repeatedly on the recogniser's guards, using
\texttt{Expr.eqv\_sort} to read the sort, closing each branch with the matching preservation
lemma.
\end{proof}

\section*{Caveats}
\label{sec:vpc-contribution}
\label{sec:vpc-review}
\begin{itemize}
\item Many reflection schemas, and every node past \ref{vpc:thm:prim-trnat-trbool}, depend on
\texttt{VEnv.WF.orderedStrong}, whose only producer is the admitted
\texttt{VEnv.WF.patsStrong}; \texttt{checkDef.WF} and every \texttt{PrimitiveResult} inherit it.
\item Nodes that chain two definitional equalities, or invert a $\Pi$
(\ref{vpc:thm:prim-bitwise-operand}, \ref{vpc:lem:prim-isdefequ-app-congr} and its consumers), go
through \texttt{IsDefEq.uniq} and \texttt{IsDefEqU.forallE\_inv}, both \texttt{sorry}-backed.
\item Every node run in the typechecker monad restores caches through
\texttt{TrExprS.weakFV\_inv}, which reaches a further \texttt{sorry} at
\texttt{TrProj.weak'\_inv}.
\item \ref{vpc:def:natbindertypes}'s syntactic Nat-binder guard is a hypothesis of the
well-founded-recursion theorems, not something the recogniser checks.
\item \ref{vpc:ind:addsconsts} and the cluster it supports through
\ref{vpc:thm:checkinductive-wf} are pre-staged: nothing constructs an \texttt{AddsConsts} yet,
and \ref{vpc:thm:checkinductive-wf} has no consumer.
\item \texttt{Nat.bitwise}'s higher-order obligation (\ref{vpc:thm:data-mkresultbitwise}) is the
only branch whose specification quantifies over every extension's operators.
\end{itemize}
